%-----------------------------------------------------------------------------
% Załącznik III
%-----------------------------------------------------------------------------

\cleardoublepage
\section*{III. Tkaczykowie z~Cisia}
\label{sec:tkaczykowie}

Najstarszym znanym autorowi niniejszej książki przedstawicielem rodziny
Tkaczyków z~Cisia jest Michał Tkacz\indeksosoby{Michał}{Tkacz}{?}{Dębe Wielkie},
urodzony około 1790~roku. Pomimo dostępu do~aktu ślubu Michała z~Wiktorią
Miros\indeksosoby[Miros]{Wiktoria}{Tkacz}{1798}{Dębe Wielkie} z~1821~roku
z~parafii pod wezwaniem Narodzenia Najświętszej Maryi Panny w~Mińsku
Mazowieckim (patrz ryc. \ref{fig:mtkacz_1821}), nie udało się ustalić ani
imion jego rodziców, ani miejsca jego urodzenia -- stąd też nie znane
są~wcześniejsze losy jego rodziny. Choć w~zdecydowanej większości aktów ślubu
z~tamtego okresu takie dane są~zazwyczaj w~pełni dostępne, to~w~przypadku
Michała Tkacza nie zostały one ujawnione -- dla jego małzonki Wiktorii Miros,
w~akcie ślubu podano zarówno imiona jej rodziców (Jacenty oraz Franciszka),
jak i~fakt, iż~urodziła się ona w~Dębem. Brak tych informacji dla Michała
wynika prawdopodobnie z~faktu, iż~nie mógł on~uzyskać metryki swojego chrztu
z~parafii, gdzie został on~ochrzczony. Jego wiek oraz wyznanie musiały zostać
potwierdzony przez lokalny, dla Dębego, sąd pokoju znajdujący się w~tamtym
czasie w~Siennicy -- świadczy o~tym wpis z~przywoływanego powyżej aktu ślubu
Michała i~Wiktorii (pisownia oryginalna):

\begin{centeredquote}
    \textit{... Stawił się Pracowity Michał Tkacz kawaler z~Dembego, podług
    złożonego przed Nami Aktu Urodzenia przez Wójta Gminy Dembe Wielkie
    odczytanego, a~przez Sąd Pokoju Powiatu Sienickiego potwierdzonego,
    lat Trzydzieści i~jeden mający, w~Dembem ze~Służby Utrzymujący się ...}
\end{centeredquote}

To~co wiemy na~pewno to~fakt, iż Michał Tkacz w~dniu swojego ślubu miał około
trzydzieści lat, od~dłuższego czasu mieszkał i~pracował w~służbie w~Dębem,
prawdopodobnie w~roli parobka. Dnia 28~stycznia 1821~roku Michał ożenił się
z~dwudziestotrzyletnią panną, Wiktorią Miros (Mirosz), która wywodziła się
z~zamieszkującej w~Dębem od~kilku pokoleń rodziny Mirosów\footnote{Wspominany
już w~niniejszej książce \enquote{Inwentarz dóbr dzierżawy Dembe}
\cite{debe1}~sporządzony w~1772~roku przez rotmistrza Andrzeja Olędzkiego
wymienia trzy odrębne rodziny Mirosów zamieszkujących wieś Dembe -- Wiktoria
prawdopodobnie wywodziła się z~jednej z~nich.}. Wpis z~aktu ślubu dotyczący
danych osobowych Wiktorii przedstawia ją~w~następujący sposób (pisownia
oryginalna):

\begin{centeredquote}
    \textit{... tudzież Wiktorya Miroszówna panna dowodząca złożoną przed Nami
    Metryką wyjętą z~księgi kościoła parafii Mińskiey iż~ma~lat dwadzieścia
    trzy, w~domu Rodzicielskim zostająca, w~Assystencyi Rodziców swoich
    Jacentego i~Franciszki Miroszów Gospodarzy z Dembego ...}
\end{centeredquote}

Wszystko wskazuje więc na~to, że~Michał Tkacz około 1810~roku przybył
do~Dębego spoza terenu mińskiej parafii, zatrudnił się tam w~roli parobka,
a~po paru latach ożenił się z~lokalną panną z~domu Miros. Nowożeńcy kilka
lat po~ślubie przeprowadzili się do~pobliskiej wsi Cisie, zakładając w~ten
sposób \textbf{rodzinę Tkaczyków z~Cisia}.

\begin{figure}[!ht]
    \centering \includegraphics[width=0.85\linewidth]{
        1821_Michał_Tkacz_Wiktoria_Mirosz_akt_ślubu_parafia_Mińsk_Mazowiecki_wpis_8.jpg}
    \captionsetup{format=hang}
    \caption{Akt ślubu Michała Tkacza i~Wiktorii Mirosz -- par. Mińsk
    Mazowiecki 1821~rok (8/1821) \cite{par_minsk0}.}
    \label{fig:mtkacz_1821}
\end{figure}

Państwo młodzi Michał i~Wiktoria Tkaczowie tuż po~ślubie zamieszkali w~Dębem,
gdzie w~1822~roku urodził się ich najstarszy syn
Jan\indeksosoby{Jan}{Tkacz}{1822}{Dębe Wielkie} (ur.~1822~r.~-- zm.~1862~r.),
po~czym około 1824~roku przeprowadzili się do~Cisia, gdzie urodziła się
pozostała siódemka ich dzieci:

\begin{itemize}
\item Marcin\indeksosoby{Marcin}{Tkacz}{1824}{Dębe Wielkie}
(ur.~1824~r.~-- zm.~?),
\item Apolonia\indeksosoby{Apolonia}{Tkacz}{1826}{Dębe Wielkie}
(ur.~1826~r.~-- zm.~1826~r.),
\item Piotr\indeksosoby{Piotr}{Tkacz}{1827}{Dębe Wielkie}
(ur.~1827~r.~-- zm.~1879~r.),
\item Kunegunda\indeksosoby{Kunegunda}{Tkacz}{1830}{Dębe Wielkie}
(ur.~1830~r.~-- zm.~?),
\item Paweł\indeksosoby{Paweł}{Tkacz}{1833}{Dębe Wielkie}
(ur.~1833~r.~-- zm.~?),
\item Łucja\indeksosoby[Tkacz]{Łucja}{Lech}{1836}{Dębe Wielkie}
(ur.~1836~r.~-- zm.~1886~r.),
\item Stanisław\indeksosoby{Stanisław}{Tkacz}{1841}{Dębe Wielkie}
(ur.~1841~r.~-- zm.~1843~r.).
\end{itemize}

Z~ósemki dzieci Michała i~Wiktorii Tkaczyków tylko trójka prawdopodobnie
dożyła wieku dorosłego -- Jan, Piotr oraz Łucja.